\section{Exploration vs. Exploitation}

\begin{frame} 
\mode<presentation>{
    \begin{center} \huge
        \secname
    \end{center}  
    \vspace{5mm}
    }
	\begin{table}[h]
	\centering
	\resizebox{0.8\textwidth}{!}{%
	\begin{tabular}{ccc}
	exploration                                                                               &  & exploitation                                                                                                         \\
	\multicolumn{3}{c}{$\xleftrightarrow{\hspace*{5cm}}$}                                                                                                                                                                               \\
	\begin{tabular}[c]{@{}c@{}}visit all state-action pairs\\ sufficiently often\end{tabular} & \hspace{2cm} & \begin{tabular}[c]{@{}c@{}}converge fast and generate\\ as much reward as possible\\ \\ (greedy policy)\end{tabular}
	\end{tabular}%
	}
\end{table}
	
\mode<presentation>{
%\begin{center}
		%\includegraphics[width=2cm]{img/explore}
%\end{center} 
%\begin{table}[]
%\centering
%\resizebox{0.2\textwidth}{!}{%
%\begin{tabular}{|c|c|c|}
%\hline
 %&                   &                       \\  \hline
 %& &  \\ \hline
 %&                       &                       \\ \hline
%\end{tabular}%
%}
%\end{table}
	%\slidesonly{\textbf{(see blackboard)}}
}
\end{frame}

\subsection{Finding a balance}

\begin{frame}\frametitle{\subsecname}


Some of the possible approaches:
\begin{enumerate}

	\item parametrize the extraction of the policy from the Q-values:
	
		%\begin{enumerate}[a\notesonly{.}]
		\begin{itemize}
		
		\item $\varepsilon$-greedy\only<2>{:

			\begin{equation}
			   \begin{array}{cc}
				 &\varepsilon\text{-greedy}\\
				 &\text{policy}
			   \end{array}
			   \pi(\vec a_k | \vec x_i) = \left\{
			   \begin{array}{ll}
				 {1-\varepsilon}, &\text{if } k = \argmax_{1 \leq l \leq A} Q(\vec x_i, \vec a_l) \\
				 {\frac{\varepsilon}
				 {A-1}}, &\text{otherwise.}
			   \end{array}
				 \right.
			\end{equation}
			
			\question{How does $\varepsilon$ control the balance?}
			}
			
		%\item 
		
		%\begin{equation}
			%\pi(\vec a_k \,|\, \vec x_i) 
				%\quad=\quad 
				%\frac{\exp\big( \beta \, Q(\vec x_i, \vec a_k)\big)}
				%{\sum_{l=1}^A \exp\big( \beta \, Q(\vec x_i, \vec a_l) \big)}
		%\end{equation}
		

		%\end{enumerate}
		\end{itemize}

	\item Optimistic Initialization
	
	\only<3>{
	A greedy sampling policy would fail at exploration,\\
	unless all Q-values are initialized with high values.
	
	\begin{itemize}
		\item initializing all Q-values with high 
			$Q_{(0)} \geq \frac{r_\text{max}}{1-\gamma}$ directs exploration
		\vspace{1mm}
		\item unexplored actions are initially made very attractive
		
		\question{How does the transition to exploration happen?}\\
		
		\notesonly{
		- Through experience, as Q-values are reduced over time exploitation takes over.
		}
	\end{itemize}
	}


\end{enumerate}

\end{frame}

